\section{문서의 독자, 범위, 상태 언어}
\label{sec:orientation}

\subsection{무엇을 읽기 쉽게 만들려는가}

이 작업공간의 기존 \sourcepath{latex/} 문서는 주차별 연구개발 기록이다.
반면 이 문서는 수학자와 암호이론가가 다음 세 질문에 바로 답할 수 있도록
정리 순서를 바꾼다.

\begin{enumerate}
  \item 논문 속 어떤 수학적 등식이 구현의 어떤 바이트 불변식으로 내려오는가?
  \item 현재 EasyCrypt가 실제로 검증한 가장 강한 정리는 무엇이며, 그 전제는
        무엇인가?
  \item 그 정리가 전체 기능정확성이나 EUF-CMA 정리가 되려면 무엇이 더
        필요한가?
\end{enumerate}

본문의 사실 우선순위는 실제 EasyCrypt 선언, \sourcepath{CLAIM_LEDGER.md},
추출 매니페스트, 최신 주차별 보고서 순이다. 주차별 문서의 앞부분은 해당
시점의 역사적 스냅샷이므로 현재 상태 판정에는 다음 자료를 우선한다.

\begin{itemize}[leftmargin=1.5em,itemsep=0.1em,topsep=0.3em]
  \raggedright
  \item \sourcepath{WEEK16_KG_NTT_MUL_}\allowbreak\sourcepath{REPORT.md},
        Week~16 KG-NTT 감사 메모 \sourcepath{latex/sections/}\allowbreak
        \sourcepath{33-week16-kg-ntt-mul-stop.tex};
  \item \sourcepath{WEEK16_KG_REPORT.md},
        \sourcepath{WEEK16_MINCORE_PLAN.md},
        \sourcepath{logs/verify-all-week16-}\allowbreak
        \sourcepath{kg-first-task.log};
  \item \sourcepath{WEEK15_REPORT.md},
        \sourcepath{RANS_ACTUAL_SUCCESS_}\allowbreak\sourcepath{WITNESS.md},
        Week~15 연구 메모 \sourcepath{latex/sections/}\allowbreak
        \sourcepath{31-week15-rans-actual-}\allowbreak
        \sourcepath{success-witness.tex};
  \item \sourcepath{WEEK14_REPORT.md}, Week~14 연구 메모
        \sourcepath{latex/sections/}\allowbreak
        \sourcepath{30-week14-hbz-full-wrapper-}\allowbreak
        \sourcepath{composition.tex},
        \sourcepath{HBZ_FULL_WRAPPER_}\allowbreak\sourcepath{COMPOSITION.md};
  \item \sourcepath{WEEK13_REPORT.md}, Week~13 연구 메모
        \sourcepath{latex/sections/}\allowbreak
        \sourcepath{29-week13-rans-core-}\allowbreak\sourcepath{composition.tex},
        독립 검증 기록 \sourcepath{logs/week13-}\allowbreak
        \sourcepath{independent-verifier.md};
  \item \sourcepath{PRODUCT_REPLAY_}\allowbreak\sourcepath{DECISION.md},
        \sourcepath{SIGNATURE_CODEC_}\allowbreak\sourcepath{OBLIGATIONS.md},
        \sourcepath{MODE2_HBZ_CODEC_}\allowbreak\sourcepath{OBLIGATIONS.md};
  \item \sourcepath{RANS_PROOF_}\allowbreak\sourcepath{DESIGN.md},
        \sourcepath{RANS_BYTE_STACK_}\allowbreak\sourcepath{INVARIANT.md},
        \sourcepath{RANS_ENCODER_}\allowbreak\sourcepath{INVARIANT.md},
        \sourcepath{RANS_DECODER_}\allowbreak\sourcepath{INVARIANT.md}와 운영 문서.
\end{itemize}

줄 번호는 편집에 따라 변하므로 본문은 파일 경로와 선언 이름을 안정적인
식별자로 쓴다.

\subsection{한글(영어) 병기 원칙}

이 문서는 수학자와 암호이론가가 구현 용어를 원문과 대응시킬 수 있도록 다음
규칙을 사용한다.

\begin{enumerate}
  \item 기술 용어가 처음 나올 때는 한글을 먼저 쓰고 영어 원어를 괄호에 둔다.
        예를 들어 정제(refinement), 불변식(invariant), 추적(trace), 프레임
        조건(frame condition), 접두부(prefix), 접미부(suffix)처럼 쓴다.
  \item 같은 문맥에서 반복할 때는 한글을 우선한다. 다만 코드와 직접 대조하는
        문장에서는 이미 병기한 영어를 다시 쓸 수 있다.
  \item EasyCrypt의 정리명, 절차명, 주장 식별자, 파일명과 코드 조각은 정확한
        검색을 위해 번역하지 않는다. \identifier{rANS}, \identifier{HBZ},
        \identifier{W32}, \identifier{W64}, \identifier{EUF-CMA} 같은 표준 약어도
        그대로 둔다.
  \item ``actual''은 실제 생성 절차를 뜻할 때 ``실제(actual)'', ``canonical''은
        ``정준(canonical)'', ``success-conditioned''는
        ``성공 조건부(success-conditioned)''로 처음 병기한다.
\end{enumerate}

\subsection{상태 용어}

\begin{description}[style=nextline,leftmargin=2.7cm]
  \item[\statusproved]
    표시된 전제 아래의 EasyCrypt 정리가 새 캐시 없이 컴파일된다. 부모 목표의
    전제가 모두 닫혔다는 뜻은 아니다.
  \item[\statuspartial]
    부모 목표에 필요한 진짜 하위 정리는 있으나, 부모 목표 전체에는 미해결
    전제나 호출 경계가 남아 있다.
  \item[\statusspecified]
    목표의 타입과 논리식은 컴파일되지만 이를 증명하는 정리는 아직 없다.
  \item[\statusblocked]
    필요한 추출, 메모리 의미론, 확률분포, 수치오차, 또는 암호학적 환원이
    구체적으로 빠져 있어 다음 합성을 진행할 수 없다.
  \item[\statusdeferred]
    필요한 간극은 정확히 식별됐지만 현재 허용된 증명 규칙과 범위에서는 닫히지
    않아 명시적 논문 경계로 동결했다. 증명되었거나 불필요하다는 뜻이 아니다.
\end{description}

\begin{boundary}[\texttt{PROVED}의 해석]
이 문서에서 \statusproved 는 항상 ``표시된 전제를 만족하는 실행에 대한 해당
정리가 검증되었다''는 뜻이다. 특히 Hoare 부분정확성 정리의 \statusproved 는
종료, losslessness, 출력분포의 일치를 함의하지 않는다.
\end{boundary}

\subsection{스냅샷과 대상}

\begin{table}[ht]
\centering
\caption{이 문서가 고정하는 검증 경계}
\label{tab:snapshot}
\begin{tabularx}{\textwidth}{L{0.27\textwidth}Y}
\toprule
항목 & 고정된 대상 \\
\midrule
논문 명세 & \sourcepath{CryptoLab_HAETAE.pdf}, 46쪽, 2026-02-04 생성,
SHA-256 \\
& \identifier{80d28b3af9af2a27dba0eb4746f9b575}\newline
\identifier{5e05cb9a86bb6596f45fda06c605f3f6}
\cite{haetae-spec} \\
파라미터 집합 & HAETAE-2만: \(n=256\), \((k,\ell)=(2,4)\),
\(q=64513\), \(\eta=1\), \(\tau=58\) \\
구현 & 고정된 \sourcepath{haetae-ref-jasmin/} 소스와
\identifier{jasmin2ec} 생성물 \cite{haetae-jasmin} \\
작성 증명 & \sourcepath{haetae-topdown-easycrypt/easycrypt/} 아래 매니페스트된
77개 EasyCrypt 대상; 마지막 보존 전체 완료 기준선은 76개 \\
최신 판정 & \identifier{STOP-KG-NTT};
\identifier{OBL-MINCORE-KEYGEN}은 실제 matrix/finalizer 스냅샷 경계에서
\statuspartial 이고 KeyGen은 KG-2/finalization에서 동결; 활성 전선은 Sign \\
재사용 증명 & 고정된 KeyGen, NTT, sampler, singularity, FFT 결과를 원래 전제
그대로 가져오는 import-only 어댑터 \cite{haetae-easycrypt} \\
제외 범위 & 다른 파라미터 집합, 물리적 부채널, fault, RNG 품질, 그리고
EUF-CMA로부터 자동으로 따라오지 않는 메모리 안전성 및 SCT \\
\bottomrule
\end{tabularx}
\end{table}

\begin{boundary}[77대상 매니페스트와 완료 증거]
현재 77번째 대상 \sourcepath{Mode2KeygenNttMulBridge.ec}는 stop-boundary 감사
결과와 actual harness의 두 active row consequence를 개별 fresh compile한다.
그 직전 76대상 전체 완료 실행은
\sourcepath{logs/verify-all-before-}\allowbreak
\sourcepath{kg-ntt-mul.log}에 보존된다. 최종 77대상 전체 완료 로그는 아직
없다. 실행 중 덮어쓰는
\sourcepath{logs/verify-all-}\allowbreak\sourcepath{summary.txt}는 마지막 줄에
\texttt{RESULT PASS authored-targets=77 cache=-no-eco}가 있을 때만 현재
77대상 전체 완료 증거다. 현재는 77번째 대상의 개별 fresh compile과 76대상
aggregate 완료만 확인된다. Hoare 정리의 \statusproved 는 여전히 실제
절차의 종료나 확률 1 실행을 뜻하지 않는다.
\end{boundary}

\begin{takeaway}{이 문서가 주장하지 않는 것}
완전한 KeyGen 분포 일치, 전체 Sign/Verify 기능정확성, 공용 문맥 추적
(public-context transcript) 일치, 승인된 원시 API 추적의 직접 \(\mu\) 등식,
전체 1474바이트 서명 코덱, 일반 정준 입력의 실제 rANS 성공 도달,
인코더/디코더 종료,
실제 전체 HBZ 래퍼의 모든 정준 입력에 대한 무조건 왕복, (h) 코덱,
full-NTT 표현과 논문 \(A_{\mathrm{gen}}s_{\mathrm{gen}}\) 곱의 동일시,
논문 KeyGen 식, MINCORE Sign/Verify 합성, 인코딩 무손실, 또는 Jasmin 구현의
종단간 EUF-CMA 보안은 아직 증명되었다고 주장하지 않는다. 고정 all-zero
입력의 성공 사후조건은 증명됐지만 그 실행의 종료/losslessness는 별도다.
\end{takeaway}
